Die Klimaanlage ist eine Kühlmaschine: Sie entzieht dem Raum ($T_\text{i} = \Ti$) Wärme und gibt sie an die wärmere Aussenluft ($T_\text{a}$) ab. Ihre Leistungsziffer kann höchstens (FoTaBe S.~172)
\begin{align}
 \eta_\text{K} &= \frac{T_\text{i}}{T_\text{a}-T_\text{i}}
\end{align}
sein. Je wärmer es draussen ist, desto kleiner wird dieser Höchstwert. Die gegebene Leistungsziffer ist also nur bis zur Aussentemperatur
\begin{align}
 T_\text{a} &= \TaF \\
   &= \Ti\cdot\left(1+\frac{1}{\eKO}\right) \\
   &= \Ta
\end{align}
möglich, das sind
\begin{align}
 \vartheta_\text{a} &= \taF = \ta \approx \result{\taP{0}{2}}
\end{align}
